\section{结论与展望}
\subsection{主要结论}
本文基于国家级时间序列数据，构建了融合EKC检验、扩展STIRPAT、中介分析、滚动一步预测与多情景模拟的综合框架，对我国碳达峰路径进行了系统研究。研究得到以下结论。

第一，煤炭占比在扩展STIRPAT估计中表现为统计上最稳健的正向驱动因素，说明在总量层面高碳能源依赖仍具较强耦合关系。

第二，人均GDP、能源强度、城镇化与产业结构等变量虽多呈正向点估计，但在2005—2022年小样本与多变量设定下置信区间较宽并跨越零，当前证据更适合支持“方向性判断”而非“强显著性断言”。

第三，技术进步变量在统计上不显著（置信区间跨越零），提示现有代理指标对减排技术扩散的刻画仍有限，未来需进一步优化技术指标体系。

第四，模型比较结果揭示了小样本约束下的“复杂度惩罚”现象。AR2 模型在滚动外推中表现最稳健，而 XGBoost、随机森林等复杂机器学习模型因严重过拟合导致误差急剧放大（MAPE 超过 14\%）。这进一步印证了预测评估应以真实外推表现为标准，在样本受限场景下盲目追求算法复杂度反而会损害预测可靠性。

第五，情景模拟表明政策强度决定峰值平台高度。低碳和深度脱碳情景均明显优于基准情景，意味着政策协同能够有效降低未来减排成本并改善峰后下降空间。

\subsection{研究贡献}
本文的贡献主要体现在三个方面。其一，在统一数据口径下将机制识别、预测比较和情景模拟贯通，形成完整证据链。其二，采用滚动一步预测检验模型外推能力，提高了预测比较的真实性和可复核性。其三，在实证结果呈现上注重图表与文字的紧密配合，使得统计结论的推导过程更具直观性

\subsection{局限与未来研究方向}
本文仍存在一定局限。首先，国家级时间序列样本长度有限，对复杂模型和严格因果识别形成约束。其次，技术进步与治理变量主要采用代理指标，尚不足以完整刻画创新质量、扩散效率与政策执行强度。再次，本文聚焦国家总量路径，未展开区域与行业异质性分解。

未来研究可沿三条路径深化：第一，延长观测样本并引入更高频、高连续性数据，以提升时序估计稳定性；第二，补充碳市场、绿色金融、能源价格和政策强度等制度变量，增强机制识别深度；第三，在国家级框架基础上向区域和行业层面递进拓展，构建“国家总量—区域差异—行业机制”联动研究体系。

在应用层面，本文框架可直接服务于达峰路径的年度复盘与政策校准。具体可将“滚动预测误差、结构变量变化率、情景偏离度”作为三类监测指标：当预测误差持续扩大时，说明模型假设与现实结构出现偏离，需要及时重估参数；当结构变量改善速度低于情景设定时，说明政策执行强度不足，应优化政策组合；当不同情景路径差距持续扩大时，可据此评估强化政策的边际收益并动态调整减排节奏。

总体来看，本文表明我国碳达峰路径优化的关键在于结构协同而非单点突破。未来政策需要在稳增长约束下持续推进能源替代、产业升级、技术扩散与治理协同，才能实现更低峰值、更稳下降和更高质量的绿色转型。
